%--- E15: Section 7.4, comment on the two approaches ------------------------
\begin{explained}{future comments: the linear and the nonlinear approach compared}
You are not identifying the \emph{functions}; you are identifying the \emph{equations or relations} that the observed functions satisfy. Both approaches share this philosophy but use different mathematical objects.

\medskip\textbf{Linear case: look for an annihilating operator.} Given an observed $u(t,x)$ one searches for a constant-coefficient operator with $P(\partial_t,\partial_x)u=0$. For $u=e^{x+t}$: $(\partial_t-1)u=0$ and $(\partial_x-1)u=0$. One studies $\Ann(u)=\{P\in\C[\partial_t,\partial_x]:Pu=0\}$ and wants to characterize all the linear operators that annihilate $u$.

\medskip\textbf{Nonlinear case: search inside a finite dictionary.} For Burgers' $u_t+uu_x=0$ the product $uu_x$ is not a linear constant-coefficient operator applied to $u$. One chooses a finite dictionary, e.g.\ $\mathcal D=\{u_t,u_x,uu_x,u_{xx},u^2\}$, and searches for coefficients, not all zero, such that $c_1u_t+c_2u_x+c_3uu_x+c_4u_{xx}+c_5u^2=0$. This is a \emph{linear-algebra} problem in the unknown coefficients even though the candidate terms are nonlinear functions of $u$. If the data satisfy Burgers the coefficient vector $(1,0,1,0,0)$ works. The dictionary can be enlarged (higher-degree products, higher-order derivatives) but its limits must be decided in advance.

\medskip\textbf{What stays and what changes.}
\begin{center}\small
\begin{tabular}{@{}L{2.9cm}L{5.6cm}L{5.8cm}@{}}
\toprule
 & Linear approach & Nonlinear approach\\\midrule
Object observed & the function $u$ & the function $u$\\
Candidate relations & linear differential operators & differential polynomials\\
Search space & polynomial ring of constant-coefficient operators & a chosen finite dictionary of jet monomials\\
Main question & which operators annihilate $u$? & which candidate relations does $u$ satisfy?\\
Main limitation & structure of the annihilator & the dictionary may omit the true equation or contain too few informative terms\\
\bottomrule
\end{tabular}
\end{center}
The common principle is \emph{identifiability}: does the observed data contain enough information to distinguish the governing equation from alternative candidates?

\medskip\textbf{Where components and interactions enter.} For $u=u_1+u_2$ one can look at the components, the total, and an interaction term such as $\Eint=u-u_0-(u_1-u_0)-(u_2-u_0)$ (the additive form of Section~\ref{sec:terms}). \emph{Decomposition and interaction analysis} ask how the solution is assembled and how its components influence one another; \emph{equation identification} asks which differential relation governs the observed solution or its components. A linear annihilator may be used for a component when the assumptions allow, and a nonlinear dictionary to identify a nonlinear equation governing the total solution: complementary analyses, not automatically the same procedure.

\medskip\textbf{Key point.} The nonlinear dictionary approach generalizes the identification philosophy; it is not a direct replacement for the annihilator in every mathematical sense. The dictionary makes the search computationally manageable, but it also limits what can be discovered.
\end{explained}

